\documentclass[a4paper,10pt]{report}\input{../0.Préambule/Préambule_cours_prof}\begin{document}

\newpage
\section*{Probabilités}
\addcontentsline{toc}{section}{Probabilités}

\vspace{-4mm}
\begin{exer1}
Pour chacun des événements suivants, indique s'il relève du hasard et si oui, place le sur l'échelle ci dessous comme dans l'exemple.

\noindent \begin{minipage}{9cm}
\begin{enumerate}[font=\color{exer1}]
\item Obtenir pile au jeu de pile ou face
\item La fête nationale aura lieu le $14$ juillet.
\item Un élève aura un tee shirt blanc demain.
\item Obtenir $6$ avec un dé à $6$ faces.
\item Trouver la bonne combinaison au loto.
\item Demain il fera beau.
\end{enumerate}
\end{minipage}
\hspace{5mm}
\begin{minipage}{8cm}
\includegraphics[scale=0.3]{../40.Image/16. Exo-Probabilité 1.png}
\end{minipage}
\end{exer1}

\vspace{-1mm}
\begin{exer1}
Une roue de loterie est partagée en huit \\ secteurs identiques numérotés de $1$ à $8$.

\noindent Gradue l'échelle ci-dessous et place-y les événements \\suivants.

\vspace{-1.5cm}
\noindent \begin{minipage}{9cm}
\vspace{1.5cm}
\begin{enumerate}[font=\color{exer1}]
\item \og Obtenir $2$ \fg{} 
\item \og Obtenir un multiple de $2$ \fg{} 
\item \og Obtenir un nombre supérieur à $4$ \fg{} 
\end{enumerate}
\end{minipage}
\hspace{5mm}
\begin{minipage}{8cm}
\includegraphics[scale=0.3]{../40.Image/16. Exo-Probabilité 2.png}
\end{minipage}
\end{exer1}

\vspace{-1mm}
\begin{exer1}
On tire au hasard un jeton parmi vingt-six jetons marqués chacun d'une lettre différente de l'alphabet. Gradue l'axe ci-dessous et place-y les probabilités des événements suivants

\begin{center}
\begin{tikzpicture}[scale=0.5]
\def\xY{-5};
\def\yY{-1};
\def\xZ{20};
\def\yZ{1};
\coordinate (Y) at (\xY,\yY);
\coordinate (Z) at (\xZ,\yZ);
\draw[step=0.5cm, lightgray, very thin] (Y) grid (Z);
\draw[step=1cm, gray ,thin] (Y) grid (Z);
\draw[ ->,line width=1.7pt,exer1] (\xY, 0) -- (\xZ, 0);
%\foreach \x in {0,1}	\draw[very thick](\x,0.2cm)--(\x,-0.2cm) node[below,fill=white]{\footnotesize \x};
\end{tikzpicture}
\end{center}

\vspace{-3mm}
\begin{enumerate}[font=\color{exer1}]
\begin{tabularx}{\linewidth}{*{3}{>{\arraybackslash}X}}
\item \og Obtenir un Z \fg{} &
\item \og Obtenir une consonne \fg{} &
\item \og Obtenir une voyelle \fg{} \\
\end{tabularx}
\end{enumerate}
\end{exer1}

\vspace{-1mm}
\begin{exer1}
On lance un dé non truqué à six faces.

\noindent Pour chacun des événements aléatoires ci-dessous, indique s'il est impossible, improbable, peu probable, probable, très probable, certain.

\vspace{-2mm}
\begin{enumerate}[font=\color{exer1}]
\begin{tabularx}{\linewidth}{*{2}{>{\arraybackslash}X}}
\item \og Obtenir un nombre inférieur a $6$ \fg{} &
\item \og Obtenir deux \fg{} \\
\item \og Obtenir un multiple de $3$ \fg{} &
\item \og Obtenir un multiple de $7$ \fg{} \\
\item \og Obtenir un diviseur de $7$ \fg{} &
\item \og Obtenir un diviseur de $60$ \fg{} \\
\end{tabularx}
\end{enumerate}
\end{exer1}

\vspace{-1mm}
\begin{exer2}
Trois personnes, Aline, Bernard et Claude, ont chacune un sac contenant des billes.

\noindent Chacune tire au hasard une bille de son sac. Le contenu des sacs est le suivant :

\medskip
\noindent \begin{minipage}{4.5cm}
\begin{center}
Sac d'Aline

\medskip
\boxed{\text{5 billes rouges}}
\end{center}
\end{minipage}
\begin{minipage}{5cm}
\begin{center}
Sac de Bernard

\medskip
\boxed{\text{10 billes rouges et 30 billes noires}}
\end{center}
\end{minipage}
\hspace{1.5cm}
\begin{minipage}{5cm}
\begin{center}
Sac de Claude

\medskip
\boxed{\text{100 billes rouges et 3 billes noires}}
\end{center}
\end{minipage}

\medskip 
\noindent Laquelle de ces trois personnes a-t-elle la plus grande probabilité de tirer une bille rouge ?
\end{exer2}

\vspace{-1mm}
\begin{exer2}
On écrit sur les faces d'un dés équilibré à six faces chacune des lettres suivantes : \textbf{NOTOUS}.

\noindent On lance le dé et on regarde la lettre inscrite sur la face supérieure.

\begin{enumerate}[font=\color{exer2}]
\item Quelles sont les issues de cette expérience ? 
\end{enumerate}
Détermine la probabilité de chacun des événements :

\noindent \begin{minipage}{7.5cm}
\begin{enumerate}[font=\color{exer2}]
\setcounter{enumi}{1}
\item $E_1$ : \og On obtient la lettre $O$ \fg{} 
\item $E_2$ est l'événement contraire de  $E_1$. 

Décris $E_2$ et calcule sa probabilités.
\item $E_3$ : \og On obtient une consonne \fg{} 
\end{enumerate}
\end{minipage}
\begin{minipage}{7.5cm}
\begin{enumerate}[font=\color{exer2}]
\setcounter{enumi}{4}
\item $E_4$ : \og On obtient une lettre du mot KIWI \fg{} 
\item $E_5$ : \og On obtient une lettre du mot CAGOUS \fg{} 
\item Gradue l'axe ci-dessous et place-y les probabilités des événements précédents.
\end{enumerate}
\end{minipage}

\begin{center}
\begin{tikzpicture}[scale=0.5]
\def\xY{-5};
\def\yY{-1};
\def\xZ{20};
\def\yZ{1};
\coordinate (Y) at (\xY,\yY);
\coordinate (Z) at (\xZ,\yZ);
\draw[step=0.5cm, lightgray, very thin] (Y) grid (Z);
\draw[step=1cm, gray ,thin] (Y) grid (Z);
\draw[ ->,line width=1.7pt,exer2] (\xY, 0) -- (\xZ, 0);
%\foreach \x in {0,1}	\draw[very thick](\x,0.2cm)--(\x,-0.2cm) node[below,fill=white]{\footnotesize \x};
\end{tikzpicture}
\end{center}
\end{exer2}

\begin{exer1}
Une urne contient $4$ boules rouges et $6$ boules vertes, toutes indiscernables au toucher. On tire une boule au hasard. Réponds par vrai ou faux.

\begin{enumerate}[font=\color{exer1}]
\item Il y a autant de chances d'avoir une boule verte qu'une boule rouge.
\item Il y a quatre chances sur $10$ d'obtenir une boule verte. 
\item Il y a $6$ chances sur $4$ d'obtenir une boule verte.
\item La probabilité de tirer une boule rouge est de $\dfrac{2}{5}$
\end{enumerate}
\end{exer1}

\begin{exer2}
On tire une carte au hasard dans un jeu de $32$ cartes. On considère les événements suivants :
\begin{enumerate}
\item[•] A : \og On obtient un roi. \fg{} 
\item[•] B : \og On obtient un as. \fg{} 
\item[•] C : \og On obtient un trèfle. \fg{} 
\end{enumerate}

\begin{enumerate}[font=\color{exer2}]
\item Les événements $A$ et $B$ sont-ils compatibles ? et les événements B et C ? Justifie.
\item Décris par une phrase sans négation l'événement contraire de l'événement C.
\item Propose un événement $D$ compatible avec l'événement $C$.
\item Détermine les probabilités des événements $A$, $B$, $C$ et $D$.
\end{enumerate}
\end{exer2}

\begin{exer2}
Un sac opaques contient des bonbons bleus, rouges ou verts, tous indiscernables au toucher.

\noindent Quand on tire un bonbon au hasard, on a deux chances sur conque de prendre un bonbon rouges et une chance sur deux de prendre un bonbon bleu.

\begin{enumerate}[font=\color{exer2}]
\item Quelle est la probabilité d'obtenir un bonbon rouge ou un bonbon bleu ?
\item Déduis-en la probabilité d'obtenir un bonbon vert.
\item Peux-tu estimer le nombre de bonbons dans le sac ?
\end{enumerate}
\end{exer2}

\end{document}